\subsection{Long-term thermal responses to landscape change}
\label{sec:long-term-cases}

Yan'an New District and Lanzhou New Area illustrate two settings where urban construction and vegetation change have occurred together. Historical Landsat imagery records the expansion of construction and changes in surrounding vegetation, while \chinatherm{} shows the evolution of leaf-on daytime thermal conditions during 2000--2024 (Figs.~\ref*{fig:yanan-long-term} and \ref*{fig:lanzhou-long-term}).

In Yan'an, thermal conditions in the construction zone progressively diverged from those of the surrounding loess hills. During 2000--2010, the area that later became the new district was still predominantly hilly terrain. Extensive land levelling was visible by 2015, followed by the development of roads and buildings (Fig.~\ref*{fig:yanan-long-term}a--f). The construction zone Y1 and surrounding hills Y2 initially followed similar temperature trajectories, which began to separate around 2012 (Fig.~\ref*{fig:yanan-long-term}j). Between 2000--2004 and 2020--2024, mean LST in Y1 increased from 29.36 to 30.30~°C, while that in Y2 decreased from 29.28 to 26.54~°C. Their temperature difference widened from 0.08 to 3.76~°C. Cooling in the surrounding hills contrasted spatially with concentrated warming patches within the construction zone (Fig.~\ref*{fig:yanan-long-term}g--i). Over the same periods, NDVI in Y2 increased from approximately 0.33 to 0.66; NDVI in Y1 declined markedly during construction and subsequently recovered to some extent (Fig.~\ref*{fig:yanan-long-term}k). Together, the temperature record and optical imagery capture the concurrent development of the new district and greening of the surrounding hills.

Lanzhou shows a different contrast between airport development and agricultural landscape change. The airport surroundings and cropland belt evolved into an interspersed landscape of construction areas and contiguous farmland, with extensive engineering works visible within the airport window in recent imagery (Fig.~\ref*{fig:lanzhou-long-term}a--f). Phase III expansion of Lanzhou Zhongchuan International Airport began in 2020, and its main works were completed in 2024, consistent with the recent changes visible in the imagery \parencite{lanzhou2024airport}. Between 2000--2004 and 2020--2024, mean LST in airport window L3 increased from 32.67 to 35.00~°C, while NDVI decreased from approximately 0.18 to 0.11. In western cropland L2, mean LST decreased from 31.45 to 29.28~°C as NDVI increased from approximately 0.24 to 0.42. The airport--cropland temperature difference widened from 1.21 to 5.72~°C, with airport warming and cropland cooling clearly expressed in the period-difference maps (Fig.~\ref*{fig:lanzhou-long-term}g--i). Mean LST in the main built-up zone L1 changed less, from 32.30 to 32.90~°C, with both warming and cooling patches within the zone.

Annual curves in both cases retain substantial year-to-year variation, while thermal differences among surface types widen as the landscapes evolve. Yan'an shows increasing separation between construction areas and surrounding greening hills, whereas Lanzhou exhibits contrasting changes between the airport and cropland. These cases illustrate the ability of \chinatherm{} to represent local thermal contrasts and their long-term evolution against a background of shared regional interannual variability.

\begin{figure}[p]
\centering
\figureorplaceholder{../figures/main/Figure10.pdf}{}{13cm}
\ReviewFloatBegin
\caption{Landscape change and long-term thermal evolution in Yan'an New District and the surrounding hills. a--f, Annual Landsat true-colour composites for 2000, 2005, 2010, 2015, 2020 and 2024. Y1 denotes the main construction zone; Y2 comprises the three surrounding hill polygons shown. g--i, Mean LST differences for 2010--2014 minus 2000--2004, 2020--2024 minus 2010--2014, and 2020--2024 minus 2000--2004, respectively. j, k, Annual \chinatherm{} LST and Landsat NDVI for Y1 and Y2.}
\label{fig:yanan-long-term}
\ReviewFloatEnd
\end{figure}

\begin{figure}[p]
\centering
\figureorplaceholder{../figures/main/Figure11.pdf}{}{13cm}
\ReviewFloatBegin
\caption{Landscape change and long-term thermal evolution across built-up, cropland and airport zones in Lanzhou New Area. a--f, Annual Landsat true-colour composites for 2000, 2005, 2010, 2015, 2020 and 2024. L1 denotes the main built-up zone, L2 the western contiguous cropland belt, and L3 the airport and surrounding window. g--i, Mean LST differences for 2010--2014 minus 2000--2004, 2020--2024 minus 2010--2014, and 2020--2024 minus 2000--2004, respectively. j, k, Annual \chinatherm{} LST and Landsat NDVI for each zone.}
\label{fig:lanzhou-long-term}
\ReviewFloatEnd
\end{figure}
\clearpage
